\chapter{Méthodologie de l'épreuve}

\noindent\textit{Bien connaître son cours ne suffit pas : réussir l'épreuve de
Physique-Chimie-Technologie du BEPC demande aussi de comprendre comment elle est
construite, ce que chaque mot de consigne exige, et comment gérer son temps. C'est l'objet
de ce chapitre.}
\vspace{8pt}

\section{Structure de l'épreuve}

\begin{propriete}[Ce qu'il faut savoir]
L'épreuve de Physique-Chimie-Technologie du BEPC dure \textbf{2 heures}, coefficient
\textbf{2}, notée sur \textbf{20 points}. Elle comporte toujours deux parties
\textbf{indépendantes et obligatoires} :
\begin{itemize}
\item \textbf{Partie A — Évaluation des ressources} (10 points), elle-même divisée en
deux volets : \textbf{a) Évaluation des savoirs} (5 points, questions de connaissances :
définitions, lois, schémas) et \textbf{b) Évaluation des savoir-faire} (5 points, courts
exercices d'application directe du cours).
\item \textbf{Partie B — Évaluation des compétences} (10 points) : une
\textbf{situation-problème unique}, ancrée dans un contexte concret (domestique,
technique, environnemental), avec plusieurs \textbf{tâches} à accomplir, qui demande de
\textbf{mobiliser plusieurs notions ensemble} pour répondre à un problème de la vie
courante.
\end{itemize}
\end{propriete}

\begin{remarque}
Les deux parties sont \textbf{indépendantes} : un échec sur une question de la Partie A
n'empêche pas de réussir la Partie B, et inversement. Ne jamais abandonner une partie à
cause d'une difficulté rencontrée dans l'autre.
\end{remarque}

\section{Lexique des consignes}

\begin{propriete}[Ce que chaque mot exige vraiment]
Le barème sanctionne autant la \textbf{qualité de la réponse au mot employé} que le
résultat numérique. Voici ce qu'attend chaque consigne :
\end{propriete}

\begin{center}
\small
\begin{tabular}{@{}p{3.1cm}p{9.2cm}@{}}
\toprule
\textbf{Consigne} & \textbf{Ce qui est attendu} \\
\midrule
Nommer / Citer & Un mot ou une expression précise, sans phrase développée. \\
Définir & La définition exacte du terme, avec le vocabulaire scientifique correct. \\
Calculer & Un résultat numérique, avec la formule utilisée et le détail du calcul. \\
Écrire l'équation & L'équation-bilan complète, avec les formules chimiques exactes. \\
Équilibrer & Ajuster les coefficients pour respecter la conservation des atomes. \\
Justifier / Expliquer & Une phrase précise expliquant \textbf{pourquoi}, en citant une loi ou propriété. \\
Schématiser / Représenter & Un schéma clair et légendé, réalisé avec soin. \\
Vérifier & Contrôler par le calcul qu'une valeur satisfait bien une condition. \\
Comparer & Mettre en relation deux grandeurs ou situations, avec une conclusion explicite. \\
Proposer & Une solution argumentée et réaliste, adaptée au contexte donné. \\
Conclure & Une phrase de synthèse répondant explicitement à la question posée. \\
\bottomrule
\end{tabular}
\end{center}

\begin{remarque}
« Calculer » attend presque toujours la \textbf{formule utilisée} avant l'application
numérique : une réponse juste sans formule visible perd des points, même si le résultat
final est correct.
\end{remarque}

\section{Gérer son temps (2 heures)}

\begin{propriete}[Répartition indicative]
\begin{itemize}
\item \textbf{5 min} : lecture complète du sujet, repérage des parties les plus faciles.
\item \textbf{25 min} : Partie A a) Évaluation des savoirs (questions de cours).
\item \textbf{25 min} : Partie A b) Évaluation des savoir-faire (exercices d'application).
\item \textbf{55 min} : Partie B (situation-problème), en lisant \textbf{deux fois}
l'énoncé avant de commencer (les données utiles sont souvent réparties sur plusieurs
paragraphes).
\item \textbf{10 min} : relecture, vérification des unités et de la présentation.
\end{itemize}
\end{propriete}

\begin{remarque}
Sur une question bloquante, ne pas s'acharner plus de $3$ à $4$ minutes : passer à la
suite, et revenir dessus s'il reste du temps. Une question non traitée fait perdre ses
points ; s'obstiner dessus fait perdre aussi ceux des questions suivantes, non lues.
\end{remarque}

\section{Rédiger et présenter sa copie}

\begin{propriete}[Bonnes pratiques]
\begin{itemize}
\item \textbf{Encadrer ou souligner} chaque résultat final, avec son \textbf{unité}.
\item Toujours écrire la \textbf{formule littérale} avant de remplacer par les valeurs
numériques (« $U=R\times I$, donc $U=\dots$ »), plutôt qu'un calcul sans formule.
\item Pour une équation chimique, vérifier systématiquement l'\textbf{équilibrage}
(compter les atomes de chaque côté) avant de la recopier au propre.
\item Ne jamais laisser une question \textbf{blanche} : même une formule correcte sans
calcul abouti peut rapporter un point.
\item Respecter \textbf{scrupuleusement} les unités demandées (A, V, W, N, mol/L\dots) :
une bonne valeur dans la mauvaise unité perd des points.
\item Dans la situation-problème (Partie B), \textbf{revenir au contexte} : conclure par
une phrase reliée à la question concrète posée, pas seulement un nombre isolé.
\end{itemize}
\end{propriete}

\section{Dix pièges fréquents à éviter}

\begin{synthese}[Les erreurs qui coûtent le plus de points]
\textbf{1.} Confondre corps simple (un seul type d'atome, ex. $\mathrm{O_2}$) et corps
composé (plusieurs types d'atomes, ex. $\mathrm{H_2O}$) (chapitre 1).\\[3pt]
\textbf{2.} Oublier de vérifier l'équilibrage d'une équation-bilan avant de la donner
comme réponse finale (chapitres 4, 5, 11).\\[3pt]
\textbf{3.} Confondre masse molaire $M$ (g/mol) et masse $m$ (g) dans la formule
$n=\dfrac mM$ (chapitre 3).\\[3pt]
\textbf{4.} Oublier que les cotes d'un dessin technique sont toujours les dimensions
\textbf{réelles}, quelle que soit l'échelle (chapitre 13).\\[3pt]
\textbf{5.} Confondre les règles du circuit série (intensité commune) et du circuit en
dérivation (tension commune) (chapitre 8).\\[3pt]
\textbf{6.} Croire qu'une tension élevée au transport signifie une intensité élevée :
c'est l'inverse, on élève $U$ pour \textbf{diminuer} $I$ et réduire les pertes (chapitre
9).\\[3pt]
\textbf{7.} Confondre puissance $P$ (en W, instantanée) et énergie $E=Pt$ (en Wh, ce qui
est facturé) (chapitre 10).\\[3pt]
\textbf{8.} Oublier que le monoxyde de carbone (CO), produit par une combustion
incomplète, est inodore et incolore mais toxique (chapitre 11).\\[3pt]
\textbf{9.} Confondre « recyclable » (peut être retraité) et « biodégradable » (se
décompose naturellement) (chapitre 12).\\[3pt]
\textbf{10.} Confondre le rôle du disjoncteur (protège le circuit) et celui du
différentiel (protège les personnes) (chapitre 16).
\end{synthese}

\begin{remarque}
Chacun de ces pièges est repris, avec un exemple corrigé, dans l'encadré « Piège
fréquent » de la fiche de révision du chapitre concerné.
\end{remarque}
